\section{Preliminaries}
\label{sec:prelim}

\subsection{Targeted influence maximization}

Let $G = (V, E)$ be a directed graph with $|V| = n$, and let $\omega_{uv} \in [0,1]$ be the
influence weight of edge $(u,v)$. A \emph{target set} $T \subseteq V$ is given, and the goal is
to choose a seed set $S \subseteq V$ with $|S| \le k$ maximising the expected number of
positively activated nodes in $T$. Write $\sigma(S)$ for that expectation.

\subsection{Threshold-dependent overexposure}

Each node $v$ carries a \emph{threshold window} $[\theta^\kappa_v, \theta^\tau_v]$. The pair is
drawn uniformly from the $2D$ simplex
\begin{equation}
  \Delta \;=\; \bigl\{(\theta^\kappa, \theta^\tau) \in [0,1]^2 : \theta^\kappa \le \theta^\tau \bigr\},
  \label{eq:simplex}
\end{equation}
that is, with density $2$ on $\Delta$. Here $\theta^\kappa_v$ is an activation threshold and
$\theta^\tau_v$ an overexposure threshold.

Write $A^{\mathrm{in}}_t(v)$ for the set of in-neighbours of $v$ that were positively activated
at some time $\le t$. Crucially, this set \emph{keeps} nodes that were later turned negative ---
their influence is assumed to persist. The \emph{exposure score} of $v$ at time $t$ is
\begin{equation}
  \delta(v,t) \;=\; \sum_{u \in A^{\mathrm{in}}_t(v)} \omega_{uv}.
  \label{eq:delta}
\end{equation}

States are $\{\text{inactive}, \text{positive}, \text{negative}\}$. At $t=0$ the seeds are
positive and remain so. At $t+1$ the exposure of the neighbours of the nodes that turned
positive at $t$ is updated, and each such node transitions by
\begin{equation}
  v \text{ becomes } \begin{cases}
    \text{positive} & \text{if } \theta^\kappa_v \le \delta(v,t+1) \le \theta^\tau_v,\\[2pt]
    \text{negative} & \text{if } \delta(v,t+1) > \theta^\tau_v,\\[2pt]
    \text{inactive} & \text{otherwise (re-evaluated when more influence arrives).}
  \end{cases}
  \label{eq:transition}
\end{equation}
Transitions are one-way: inactive $\to$ positive $\to$ negative, never backwards. A negatively
active node does not promote. We call this the \emph{threshold-window} process.

\begin{remark}[Activation probability]
\label{rem:lemma1}
Integrating over the window draw gives the marginal positive-activation probability as a
function of exposure alone,
\begin{equation}
  \Pr\bigl[\theta^\kappa \le \delta \le \theta^\tau\bigr] \;=\; 2\,\delta\,(1-\delta),
  \label{eq:g}
\end{equation}
which we denote $g(\delta)$. This is the quantity the source model builds its analysis on. It is
$0$ at $\delta \in \{0,1\}$ and peaks at $\delta = 1/2$.
\end{remark}

Two structural facts drive everything that follows and are worth isolating.

\begin{observation}[Exposure is monotone in time, the probability is not]
\label{obs:monotone}
$A^{\mathrm{in}}_t(v)$ grows with $t$ and never shrinks, because a node that has been positively
activated keeps contributing even after turning negative. Hence $\delta(v,t)$ is
non-decreasing in $t$ for every $v$. The activation probability $g(\delta) = 2\delta(1-\delta)$ is
therefore \emph{non-monotone in time}: it rises while $\delta < 1/2$ and falls once
$\delta > 1/2$. The only source of non-monotonicity in the model is the shape of $g$, not the
exposure accumulation.
\end{observation}

\begin{observation}[Exposure is a set function of the seed set]
\label{obs:setfunction}
Condition on a fixed draw of all windows. Then for each $v$ the indicator of ``$v$ was ever
positively activated'' is a deterministic function of $S$, and consequently so is
$\delta(v,t)$ for every $v$ and $t$, and so is the final spread. The only randomness in the
entire process is the single window draw.
\end{observation}

Observation~\ref{obs:setfunction} is the reason the usual sampling machinery has nothing to work
with, and we make that precise next.

\subsection{Reverse-reachability sampling, and what it requires}

RR sampling relies on the identity
\begin{equation}
  \mathbb{E}[\sigma(S)] \;=\; n \cdot \Pr\bigl[S \cap \mathcal{R} \neq \emptyset\bigr],
  \label{eq:rr-identity}
\end{equation}
where $\mathcal{R}$ is a random set generated by a reverse walk from a uniformly random node.
The identity is available whenever the process admits a \emph{randomised graph structure} over
edges --- the live-edge graph in independent cascade, the trigger set in linear thresholds ---
such that $v$ is active iff it is reachable from $S$ in that structure. Two features are
required: (i) reachability can start from a base case in which a node is active under
\emph{incoming influence from a single edge or a single predecessor}, and (ii) the structure is
sampled independently of $S$.
